\documentclass[conference]{IEEEtran}
%\IEEEoverridecommandlockouts
% The preceding line is only needed to identify funding in the first footnote. If that is unneeded, please comment it out.
%Template version as of 6/27/2024

\usepackage{cite}
\usepackage{amsmath,amssymb,amsfonts}
\usepackage{algorithmic}
\usepackage{graphicx}
\usepackage{textcomp}
\usepackage{xcolor}
\usepackage{multirow}
\def\BibTeX{{\rm B\kern-.05em{\sc i\kern-.025em b}\kern-.08em
    T\kern-.1667em\lower.7ex\hbox{E}\kern-.125emX}}

\begin{document}
\bibliographystyle{plain}

\title{Measuring the Impact of Monitoring on HPC Workloads}

\author{\IEEEauthorblockN{Steven Presser-Valderrama}
\IEEEauthorblockA{\textit{HPC Nexus Laborartory} \\
\textit{Technological University Dublin}\\
Dublin, Ireland \\
ORCID: 0000-0002-0362-6741}
\and
\IEEEauthorblockN{Tania Malik}
\IEEEauthorblockA{\textit{HPC Nexus Laborartory} \\
\textit{Technological University Dublin}\\
Dublin, Ireland \\
email address or ORCID}
}

\maketitle

\begin{abstract}
HPC constantly balances performance and monitoring. Significant previous work 
has examined the impact of monitoring on HPC workloads. However, the 
methodology, statistical robustness and even units used in previous work vary 
greatly. This work proposes a method for comparably measuring the impact of 
monitoring on HPC workloads. Additionally, the method is applied to measure 
the impact of popular and emerging in- and out-of band monitoring tools. 
Finally, a sample implementation of the methodology is contributed for future 
use.
\end{abstract}

\begin{IEEEkeywords}
High-Performance Computing, Monitoring, Metrics, Performance
\end{IEEEkeywords}

\section{Introduction} \label{introduction}
HPC is the art of extracting as much performance as possible from as little 
computing infrastructure as possible. This requires balancing many competing 
needs, such as processor use, memory use, cache use, and network bandwidth.
Performing this tuning requires monitoring the systems on which an HPC workload
runs. However, much like Shrödinger's infamous cat
\cite{schroedingerGegenwaertigeSituationQuantenmechanik1935}, the act of 
monitoring and
measuring HPC system performance impacts the system. This can be through the 
use of CPU time that the workload could otherwise use, jitter on the network 
caused by monitoring processes or even cache invalidations caused by
monitoring processes.

This work seeks to quantify the impact of monitoring processes on HPC workloads
in a standardized way. 
Further, this paper seeks to create a standardized method for evaluating the 
impact of monitoring on an HPC system, as well as contributing a reference
implementation thereof.

In Section \ref{background}, we discuss prior work, including evaluations of monitoring 
systems. Section \ref{motivation} discusses our motivation for undertaking this 
work. Section \ref{methodology} continues by discussing the proposed methodology to measure monitoring 
impact, with a reference implementation in Section \ref{implementation}. We continue in Section \ref{results} by 
showing the results of using the methodology, which are discussed in Section \ref{discussion}.
while Section \ref{futurework} describes future work. Finally, we conclude in Section \ref{conclusion}.

This work contributes a step-by-step framework for measuring changes caused by 
monitoring on HPC systems, standards for reproducability when using the 
framework, and a reference implementation of the framework. Finally, it applies 
the framework to provide statistically-meaningful measurements of overhead for 
HPCPerfStats, LDMS, and BMC-based out-of-band monitoring.

\section{Background} \label{background}

This work focuses on the collection of metrics in HPC systems, as well as 
methods for evaluating their impacts on the system.

A metric is a numeric measurement of a system property collected at regular intervals.
Some examples of
metrics are temperature of a core, memory usage of a system, cache misses in a
specified period, and lifetime power usage of a system. Measurements can
then be broken down into two subcategories: instantaneous (e.g.: temperature of a core) and
accumulating (e.g.: power usage or cache misses). Accumulating metrics then have 
two subcategories: windowed (e.g.: cache misses in the last minute) and 
lifetime (e.g.: lifetime power usage).

The methods used to gather metrics can be divided into in-band and out-of-band.
Collecting metrics in band is done via software running on the same computing 
hardware running the HPC workload. That is, via software running on the compute 
node's operating system. In contrast, out-of-band metrics are gathered via any 
method that does not run software on the compute node's operating system.

\subsection{In-Band Metric Collection}

Classically, metrics about how compute nodes in an HPC system are performing 
have been collected using software running on the operating system of the node. 
Because the monitoring software runs on the same processor and uses the same 
resources as the HPC workload, 
poorly-written monitoring code can cause overhead on HPC nodes, thus slowing 
down the HPC workload. For example,
this can occur if a monitoring process uses significant network bandwidth, 
which slows communications between nodes. Further, if monitoring does not 
run at the same time on all nodes involved in a workload, it can cause jitter - 
some processes are delayed and therefore slower to communicate results, causing 
all processes to wait for the slow results.
Therefore, significant effort has gone into developing low-impact/low-overhead
monitoring software that can run alongside compute workloads.

Such software can either be deployed per-job by the user or system-wide by the 
administrator. Broadly speaking, both types serve 
the same purpose: collecting performance counters accessible to the operating 
system (such as CPU frequency and memory utilization) and conveying them to a 
central data store with as little impact on the compute node and network fabric 
as possible.

Two examples of per-job metric-gathering software are LIKWID\cite{gruberLIKWID2024}
and PAPI\cite{mucciPAPIPortableInterface}. These tools are designed to be 
deployed by the user launching an HPC job and can be tuned by the user to 
collect only the required data. However, the use of these tools is not without 
cost.
Voevodin et al.\cite{voevodinOverheadAnalysisPerformance2022} found a measurable 
overhead for both LIKWID (2.78\%) and PAPI (3.83\%).

System-wide monitoring is deployed by the adminsitrator of the system. In 
contrast to per-user metric-gathering, it cannot be disabled for specific jobs. 
Uses may include hardware health monitoring, post-mortem of failed jobs, or 
classification of workload types to plan for the next system. Popular examples 
are LDMS\cite{agelastosLightweightDistributedMetric2014} and HPCPerfStats
(formerly TACCStats)
\cite{evansComprehensiveResourceUse2014,TACCHPCPerfStats2026}. 
 LDMS claims to have no statistically significant impact on several 
benchmarks\cite{LightweightDistributedMetric}, as tested on Blue Waters, a 6912 
node system. Additionally, measurement of CPU time was performed and found that 
LDMS used 0.009\% of CPU time while collecting 1012 metrics every 5 seconds
\cite{brandtHighFidelityData}. HPCPerfStats estimates an overhead of about 3\% 
with collection at 1 Hertz\cite{TACCHPCPerfStats2026}.

\subsection{Out-of-Band Metric Collection}

In contrast to in-band metrics, out-of-band matrics are gathered from any source
that is not the operating system of the node. Some examples include node 
Baseboard Management Controllers (BMCs)
\cite{liMonSTerOutoftheBoxMonitoring2020,cockEnzianOpenGeneral2022a} and 
non-contact\cite{kamilPowerEfficiencyHigh2008} and 
contact\cite{asgherEvaluatingHardwareSoftware2025,ilschePowerMeasurementTechniques2019} external power 
meters. Additionally, significant research has been performed adding internal power 
measurement devices\cite{ilschePowerMeasurementsCompute2015,larosPowerInsightCommodityPower2013,vlugtPowerSensor3FastAccurate2025}.

External measurement devices provide an easy method to acquire 
high-frequency measurements of the power used by a system. This data can be 
used for many purposes, including 
digital twins\cite{brewerDigitalTwinFramework2024}
 or to classify workloads\cite{kohlerRecognizingHPCWorkloads2021}.
However, this data is limited only to what can be measured externally from the 
node. Additionally, any monitor which requires modifying the compute node 
itself will most likely void the warranty on the hardware, making it difficult 
or impossible to recieve support or replacement parts from the manufacturer.

In contrast, the BMC is a common element of server-class hardware, designed and 
integrated by the manufacturer.
BMCs typically manage the underlying hardware and firmware of a node, provide remote power 
control capabilities, and remote console connections. Communication with BMCs 
is most often via the Intelligent Platform Management Protocol 
(IPMI)\cite{IPMIStandard}, RedFish\cite{RedfishSpecification}, 
or a web interface built by the manufacturer of the BMC
\cite{enterpriseLoggingILOWeb,PowerEdgeSetIDRAC10,HPEAgentlessManagement}.

In contrast to exeternal meters, BMCs may interact with the compute node itself. 
Often this takes the form of reading sensors via hardware busses, but some implementations may 
interact with kernel modules or the operating system\cite{IDRACServiceModule}. 
Interestingly in an HPC context, this interaction is sometimes inverted, with 
the BMC providing sensor readings to the operating system\cite{enterpriseUserAccessCompute,martinCrayXC30Power}.

However, BMCs are often significantly more limited in the speed with which they 
can collect data. For example, Lai\cite{laiBenchmarkingBMCs} tested RedFish 
response times on two BMCs and found that the time to respond to a single 
(HTTP) RedFish request varied between 0.5 seconds and 0.93 seconds. As each 
RedFish request can only fetch a sinlge metric, this severly limits the speed 
at which metrics may be collected.

Minimal previous work exists on the overhead of out-of-band monitoring on HPC 
workloads. Fasse et al.\cite{faasseOutBandPerformance2020} run 
the SPEC CPU 2107 benchmark once without monitoring and once with monitoring. 
They then examine the two results, and 
assert that because the difference (0.5\%) is with run-to-run variation of the 
benchmark, there is no performance impact. However, they do not collect a 
statistically significant sample, nor perform any statistical analysis.

Other work, such as Li et al.
\cite{liMonSTerOutoftheBoxMonitoring2020}
focuses more on how the use of a BMC for metric collection shifts the network 
traffic associated with metric collection to the management network, which 
they assert should reduce load on the compute fabric. They do not examine 
whether or not using the BMC for monitoring has any impact on the performance 
of the compute node.

In short, the general consesus is that monitoring via the BMC will have no 
effect on the running workload and that monitoring via the BMC may have a 
positive effect compared to performing monitoring via the operating system.

Out-of-band monitoring data has already found significant use in the simulation 
of HPC systems and their supporting infrastructure. For example, out-of-band 
monitoring data was critial in building 
ExaDigiT\cite{brewerDigitalTwinFramework2024}, which provides simulations of 
liquid-cooled HPC systems. Software like ExaDigiT will be critical in planning
and optimizing future HPC systems, including reducing total energy-to-solution 
in large systems.

\subsection{Overhead}

Overhead is the slowdown an HPC job experiences due to a specific configuration. 
For example, we might speak of the overhead experienced when running in 
containers\cite{rezendeallesAssessingComputationCommunication2018} or while 
running monitoring software\cite{boehmeCaliperPerformanceIntrospection2016}. 
It is generally agreed that overhead is measured by how much slower an HPC 
workload runs in comparison to some base configuration, expressed in percent.
There are generally three critical elements to reproducing 
measurements of overhead: a precise definition of the base (unaltered) case, a 
precise definition of the experimental case, and a definition of the 
statistical measurements performed to compare the cases.

However, within the subfield of HPC monitoring systems, there appears to be 
little agreement on methodology for measuring and reporting on overhead.

For instance, Voevodin et al.\cite{voevodinOverheadAnalysisPerformance2022} 
measure overhead for PAPI and LIKWID under various performance counter 
multiplexing conditions. They clearly define the base case as running DIMMon 
without 
PAPI or LIKWID on Lomonosov-2, then measure the increase in runtime when PAPI 
and LIKWID multiplexing are enabled. They do not describe how often DIMMon 
collected data, though they did keep this constant during experiments. Their 
measurement of statistical significance is "overhead is not significant [when] 
the upper bound of the confidence interval for default DIMMon is greater than 
the lower bound of the confidence interval" in the experimental case. However, 
they do not describe what confidence interval was used. Additionaly, they 
remove outliers "showing significantly different execution time" from their 
data set, but do not define significantly different in this context.
They find that 95 of their 264 tests lack statistical significance as well as 
concluding that "external influence of the HPC system can be much more 
noticeable" than the effect of the monitoring system.

HPCPerfStats\cite{TACCHPCPerfStats2026} is a system-wide in-band HPC monitoring 
system. They provide two different measurements of overhead: approximately 3\% 
with collection running at 1 Hertz\cite{TACCHPCPerfStats2026} and less than 
0.1\% with collection once every 10 minutes
\cite{evansComprehensiveResourceUse2014}.
Neither of these measurements gives a methodology (though a reader may 
reasonably assume the base case is no monitoring and the experimental is 
HPCPerfStats running collection at the described frequency).
Finally, neither measurement has any associated statistical methodology.

Even when doing experiments where proving (or disproving) an effect is critical, 
flawed methodology is used or descriptions of methodology are lacking. For 
instance, Faasse et al. \cite{faasseOutBandPerformance2020} explicitly seeks to examine 
whether or not performing monitoring via the BMC has an effect on the running 
workload. They call this "do no harm" testing. The experimental setup is well
described, including two different types of benchmarks, each of which is run 
once without BMC-based monitoring (the base case) and once with. However, they 
fail to gather any kind of statistically meaningful data. Instead, they run 
each benchmark once under each condition, note that the performance difference 
is less than 0.5\%, and end their analysis by saying the difference is
"within the run-to-run variation of the benchmark". One set of performance 
metrics with no repetition does not data make - it barely qualifies as an 
ancedote and is insuffucient proof that no harm has occurred.

In short, significant prior work testing the overhead of monitoring fails to 
adequitely describe experimental conditions for reproduction or fails to 
perform meaningful stastitcal testing on results, which makes it difficult for 
systems administrators to understand how using these systems may increase 
overhead on their systems.

\section{Motivation} \label{motivation}

System-wide monitoring is of exreme importance to operating HPC systems. The data 
is needed for everything from monitoring hardware for failure to determining the 
workloads running on a system. However, collecting that data also has a 
cost in the form of overhead - lost computation on the HPC system.

In this work, we aim to provide 
a logical and software frameowrk for performing comparisons of different 
monitoring options. Doing so will ease the selection for systems administrators 
by giving them an important insight into the overhead of each and therefore the 
tradeoffs.

We will then apply this framework to create directly comparible measurements of 
several in-band monitoring solutions, as well as to provide the first 
statistically meaningful data on whether or not monitoring via the BMC affects 
running workloads.

% This needs to go in the direction of more cabale BMCs and BMCs which may start to interact with the OS. Need to be able to verify that there is no impact.
% Further, some 
% manufactuers have BMCs which can interact with the host operating system via 
% kernel drivers (for example, in Dell's iDRAC\cite{IDRACServiceModule}) Even on 
% systems which don't use kernel drivers, 
% communication from the BMC to sensors could be via a shared bus (i.e., i2c), 
% which could have effects on the system.

% Out-of-Band monitoring is an underexplored area of 
% Previous work focused on using simulated out-of-band metrics
% for classification of jobs by type\cite{presserRealTimeClassification2022}. In
% preparing to extend that work to non-simulated out-of-band monitoring data, we
% noticed that many in-band metric systems had associated measurements of impact 
% on HPC workloads. However, we noticed that these were not standardized or 
% comparable. Further, we noted that no measurement had ever been made of the
% impact of out-of-band monitoring, at least for BMC-based monitoring.

% Therefore, our motivation is to fill two gaps in existing knowledge: first, 
% to provide directly comparable measurements of various monitoring systems, 
% including both in- and out-of-band. Second to demonstrate what (if any) impact 
% BMC-based out-of-band monitoring has on HPC workloads.

\section{Framework} \label{methodology}

We present a framework designed to determine if there is a meaningful difference between 
a base case (no monitoring) and various experimental conditions. The framework 
has two major components: a recommended experimental flow and a reporting 
framework. Together these create a known methodology that should allow a degree 
of comparison between various configurations, including between sites.

\subsection{Experimental Flow}

The general
experimental flow is:

\begin{enumerate}
\item Create a hypothesis to test
\item Determine the benchmarks to run
\item Tune benchmarks to fully utilize the system and prepare test cases
\item Perform a pilot test
\item Gather baseline data without monitoring
\item Gather experimental data with monitoring (in various configurations, if desired)
\item Perform statistical tests comparing each experimental scenario to the 
baseline scenario
\item Report the Results
\end{enumerate}

These steps are discussed in more depth in the following subsections.

\subsubsection{Create a Hypothesis}

Any good experiment begins with a hypothesis. This defines what we wish to test.
This framework is best at testing hypothesies such as "Implementing an in-band 
monitoring system will have a measurable effect on system performance" or 
"Changing option X on the monitoring system will reduce system overhead" or 
"The new version of the monitoring system will reduce system overhead".

Additionally, a threshold is selected for equivalence - how much difference 
between test cases is acceptable to still consider them to be the same. In the 
context of monitoring, this number is how much effect on the system is 
considered negligible and can be accepted.

We strongly recommend writing down the hypothesis to test before beginning 
other experimental work.

\subsubsection{Selecting Benchmarks}

We begin by selecting benchmarks to run. There is no single correct set of 
benchmarks. Instead, these will vary by system. As guidance, we suggest that 
these benchmarks include at least those used in system acceptance. Additionally, 
we suggest using any benchmarks intended to be small scale simulations of 
workloads the system is expected to run (so-called microbenchmarks).

We anticipate HPL\cite{dongarraLINPACKBenchmarkPresent2003} and 
HPCG\cite{herouxNewMetricRanking2013} being common choices of benchmark, if 
only because they are so often used in system acceptances.

\subsubsection{Tuning Benchmarks and Preparing Test Cases}

The benchmarks are then tuned to the experimental system. Specifically, each 
benchmark is tuned to fully saturate the resources it is to test. For example, a 
CPU benchmark is tuned to utilize the CPU at 100\%, an IO benchmark to saturate 
the IO system, etc. Benchmarks which simulate an application are tuned for best 
performance. We anticipate that in many cases, such tuning will already have 
been perfomed as part of system acceptance.

Additionally, test cases are prepared which reflect the experimental 
condition(s). For example, if testing whether a change to the monitoring system 
affects overhead, two nearly identical images might be prepared: one with 
the monitoring system and one without.

\subsubsection{Pilot Test}

Next, we must determine the sample size required for statistically significant 
results. A small pilot study of the benchmarks is run. This includes running 
each of the tests under the baseline and experimental conditions a 
pre-determined number of times. More runs here may result in needing fewer 
later, but there are no set rules for determining how many to run.

The gathered data is then used to perform power analysis, which determines how 
many runs are required to determine statistical significane. This number should 
be taken as a minimum - more data will allow more confidence in the results.

\subsubsection{Gathering Data}

In order to gather statistically significant data, the benchmarks are run 
for at least as many repitions as required by the power anaylsis and a
performance metric gathered from each run. The performance 
metric need not be the same between all benchmarks. Example performance metrics 
include runtime, FLOPs, and I/O operations per second, but will vary based on 
the benchmark.

Benchmarks are then run a pre-determined number of times under each 
experimental condition and the performance metric from each stored.

% Gathering data is performed in two stages. First, data is gathered with no 
% monitoring enabled. Then, the same procedure is repeated with monitoring 
% running.

% Gathering data involves running the benchmarks and collecting performance 
% metrics from the benchmarks until a statistically-significant sample is acquired.

\subsubsection{Analysis}

Once sufficient data is collected, the percent change from the baseline is 
calculated for each experimental condition. The significance of the result is 
then measured by performing two 
statistical tests: first, a T-test, to determine if there is a statistically 
significant change versus the baseline. Second, the Two One-Sided T-tests (TOST) 
test is performed to determine if the variance is within the pre-determined 
bounds.

\subsection{Reporting Results}

Comparison within a site is easy enough. However, in order to provide utility, 
this methodology must also provide comparison between sites. In order to do so,
users of the methodology shall report information on the configuration used to 
run the benchmark. We provide four tiers of reporting, ranging from the 
minimal reporting required to publish through full reproducability. These are:

\begin{itemize}
    \item Bronze: Minimal information required to report a result
    \item Silver: Results somehat comparible between sites, however may not be replicable
    \item Gold: Results are replicable, but setting up a reproduction environment may be difficult
    \item Platinum: Full reproducability. Any reader with sufficient resources can reproduce the results exactly
\end{itemize}

In the follwoing sections, we discuss the exact informaiton to be included in 
each tier. Each tier builds on the previous and must report the previous tier's 
information as well. Users of the framework are always welcome to report any 
additional information they believe is relevent. Above the bronze level, the 
required data is large enough that the authors recommend including it in an 
appendix or as associated artifacts, rather than the main body of a publication.

\subsubsection{Bronze}

At the bronze level, the reporting contains vital information about the results 
and the system. At a minimum these include:

\begin{itemize}
    \item A description of the base case and a description of how each 
    experimental case varies from the base case
    \item A list of the benchmarks run, including a description of the 
    configuration.
    \item A summary of the results, including the measured overhead, standard 
    deviation of the results, statistical methods used (if they vary from those 
    described here) and measures of statistical significance.
    \item A description of the hardware the results were obtained from, 
    including hardware manufacuturer, CPU count and model number, RAM speed and 
    amount, storage system used, network system used and other vital 
    information about the system.
    \item The name and version number of the job manager
    \item The name and version number of the operating system used
\end{itemize}

Specifically, the measured overhead should be expressed in percent overhead per 
second when collecting at 1 Hertz. Ideally, collection is performed at 1 Hertz
and transmitted to collection and aggregation servers at the same rate. 
Results should specifically note when this is not the case.

We anticipate that these will be provided in the body of any work utilizing 
this framework and have delibereately kept the reported information relatively 
small in order to enable such reporting.

\subsubsection{Silver}

In contrast, the Silver level begins to add information we anticipate being 
included in works in either appendixes or artifact repositories. It contains 
the most critical artifacts and outputs needed to verify results. Specifically, 
reporting at the Silver level adds:

\begin{itemize}
    \item Performance measurements from each benchmark run
    \item Any software used to perform statistical analysis
    \item Configuration files, if any, for the benchmarks
    \item List of services running on the compute nodes (e.g.: from 
    \texttt{systemctl status})
\end{itemize}

\subsubsection{Gold}

The Gold level adds information that should allow close reproduction
the environment that the tests were run in. Gold-level reporting adds:

\begin{itemize}
    \item A list of software versions installed on the system where the 
    benchmarks run. This could be collected with for example
    \texttt{dnf list installed}, \texttt{yum list installed}, or 
    \texttt{apt list --installed}, depending on operating system. We also 
    recommend including \texttt{/etc/os-release}, where available.
    \item Network system specification, including version numbers.
    \item Storage system specification, including version numbers.
\end{itemize}

These items are kept separate from the Silver level as some may require 
administrative access to the system. Requiring them at the Silver level would 
have forced some users of the framework to only report at the Bronze level, 
while separating these requirements to the Gold level enables all users to 
report to the Silver level, as long as they have any type of access to the 
system.

\subsubsection{Platinum}

Platnum is the highest level of reproducability and attempts to include all 
information required for an exact reproduction environment. To achieve a 
platinum level of reproducability, artifacts should include:

\begin{itemize}
    \item Full scripts and code used to run tests and statistics (excluding 
    benchmark code when publically available or prevented by the license of the 
    benchmark)
    \item Configuration files for the workload manager
    \item A list of changed files on the system, as well as the contents thereof 
    or diffs of the changed content, except when sharing the file contents would 
    be a security issue or the change was made by software, rather than an 
    administrator. The list of changes could be collected using (for example)
    \texttt{rpm -Va}. Alternatively, an image of the compute node or other 
    reproducable method of creating a compute node image could be provided.
    \item Configuration files or details of all changed options for the network
    \item Configuration files or details of all changed options for the storage 
    system.
\end{itemize}

Platinum reproducability should allow complete and exact reproduction of the 
experiment and validation of the results.

\section{Implementation} \label{implementation}

The framework described in the previous section is useful, but may be complex 
to implement. In order to ease the use of this framework for others, we provide 
\cite{} a reference implementation. This reference uses 
reFrame\cite{karakasisRegressionFrameworkChecking} to manage running benchmarks,
then performs analysis in a Jupyter notebook.

reFrame was selected as a base software because it can easily be installed and 
run on most systems without requiring administrative access. It is 
distributed as a Python package, which may be installed in a virtual environment 
and therefore does not require administrative access to install. Further, it 
does not require alteration of existing job managers to use. Instead, it is 
configured to use the user's existing access and therefore may be run without 
the assistance of the systems administrators. Finally, reFrame brings several 
convienent functions, including automatic compilation of test code and storage 
of results.

In order to validate this framework, we run it with three experimental 
conditions:

\begin{itemize}
    \item Using LDMS
    \item Using HPCPerfStats
    \item Collecting metrics from the BMC with a custom script
\end{itemize}

The benchmarks run under Slurm on a HPE DL360 Gen 10 with two Intel Xeon Gold 
5120 processors and 16 GiB DDR4 2133MHz RAM. The system has an iLO5 as a BMC.
The network was 1 Gib Ethernet, provided by the built-in network interface using 
a Broadcom NetXtreme BCM5719. Storage was provided by the built-in RAID 
controller (based on an Intel C620) and an Intel SSDSA2BW16 SSD. This node was 
the only worker in a Slurm 26.05.2 cluster and ran Rocky Linux 9.8. It was 
connected to a KVM virtual machine running the Slurm controller and NFS home 
filesystem via an unmanaged hub.

\subsection{Experimental Hypothesis}

We begin by stating our hypothesies:

\begin{enumerate}
    \item The use of in-band monitoring systems will increase overhead in an HPC 
    system
    \item The use of out-of-band monitoring will not increase overhead in an HPC
    system
\end{enumerate}

\subsection{Selecting Benchmarks}

In order to keep this evaluation simple, we focus on single-node CPU 
benchmarks. Because the software monitoring systems must use the CPU, we 
believe this is the element of the system most likely to show an effect from 
monitoring software. Additionally, we believe that if the use of a BMC for 
monitoring has any effect, it will affect the CPU, which controls almost all 
elements of the compute node.

We selected the EP (Embarassingly Parallel) benchmark from the NPB 
suite\cite{NASParallelBenchmarksa} as our test benchmark. This benchmark is 
CPU-focused and provides an "estimate of the
upper achievable limits for floating point performance"\cite{NASParallelBenchmarks}
and includes minimal interprocess communication. We believe that if the CPU 
must perform additional work for the monitoring system, that will appear as a 
slowdown or reduction in performance of the benchmark.

We selected the D and E problem sizes of the EP benchmark. The D size runs 
quickly (about 50 to 60 seconds on our test hardware), while the E takes 
approximately 15 minutes.

\subsection{Tuning Benchmarks and Preparing Test Cases}

Performance tuning in this case was extremely simple: the EP benchmark was set 
to run 56 tasks, equal to the number of hardware threads available to the operating 
system. The system was thus kept at 100\% utilization. Finally, we select 
millions of operations per second (Mops/s) as our performance benchmark. This 
value is automatically calculated by the EP benchmark.

In order to test the effect of in-band monitoring, we evaluate two different 
monitoring systems: HPCPerfStats and LDMS. The configuration used for each is
available in the artifact appendix of this paper. HPCPerfStats was configured 
with the default collection, while LDMS was configured to collect using the 
\texttt{meminfo}, \texttt{coretemp}, \texttt{loadavg}, \texttt{procnet}, 
\texttt{procstat2}, and \texttt{vmstat} modules. Both were configured to collect 
and send data to a central server once per second (1 Hertz).

In order to test the effect of BMC-based out-of-band monitoring, we use the BMC 
built into the compute node, an HPE iLO5. Because this is part of the hardware, 
we cannot test any alternatives without fundamentally changing the hardware.
The HPE iLO5 provides the ability to measure CPU, I/O and memory bus utilization 
via \texttt{amsd} (Agentless Management Service daemon)\cite{HPEAgentlessManagement},
a daemon which is installed on the operating system of the node.
Unfortunately, it appears that on the system used, \texttt{amsd} only provides 
this information once every 20 seconds. We nonetheless install it in order to 
maximize the potential impact of monitoring via the BMC.

BMC sensor data was collected from the BMC once per second using RedFish.
This collection includes the \texttt{amsd} data, although it changes only once 
every 20 seconds. The 
script and URLs used for collection are available in the artifact appendix of 
this paper. The collector ran on the Slurm admin node.

Finally, we script shifting between the four test cases (no monitoring, 
HPCPerfStats, LDMS, and BMC-based) in order to ensure consistent system state.
This script is available in teh artifact appendix of this paper.

% The following benchmarks were set to run using XX tasks, thus using all cores 
% of the system. Where tasks must be mapped to a three dimentional arrangement, 
% we have noted the X, Y, and Z task counts in parentheses.

% \begin{itemize}
%     \item 
% \end{itemize}

% The remaining benchmarks have specific requirements for the number of tasks 
% (such as a square number). In 
% order to utilize 100\% of the CPU, we disabled the unused processor cores using 
% \texttt{/sys/devices/system/cpu/cpu*/offline}\cite{KernelorgDocDocumentation}.
% The benchmarks therefore ran witht he following task counts:

% \begin{itemize}
%     \item 
% \end{itemize}

\subsection{Pilot Test}

Once the benchmarks are fully tuned, we run our pilot test. For the code used 
to run these tests, please see the artifact appendix of this paper.

Initially, we run 20 repetitions of each benchmark under the baseline 
(no monitoring) and each of the experimental conditions. We then perform a 
power analysis using the following values to determine the minimum number of 
runs of each benchmark are required. The Jupyter notebook used to perform power 
analysis is included in the artifact appendix.

The following parameters are used in the power analsysis:

\begin{table}[h!]
\centering
\begin{tabular}{| c | c |} 
    \hline
    Parameter & Value \\ [0.5ex] 
    \hline
    Repeats, each benchmark & 20 \\ 
    \hline
    $\alpha$ (alpha) & 0.05 \\ 
    \hline
    Power & 0.8 \\
    \hline
\end{tabular}
\caption{Parameters used for power analysis}
\label{table:power_params}
\end{table}

These parameters indicate that if a effect exists, there is an 80\% chance it 
will be found. Additionally, setting $\alpha$ to 5\% indicates we will tolerate a 
maximum of 5\% chance that any result is a type I error (false positive).

\subsection{Gathering Data}

Power analysis indicated that EP in size D should be run approximately 200 
times. In size E it indicated 110 repetitions were required, except in the 
HPCPerfStats case, where the 20 repetitions from the pilot study had already 
achieved statistical signficance. Therefore, no additional tests of EP in the 
HPCPerfStats case were run.

We then ran each test the number of times indicated. In the case of EP with 
size E, we ran 5 tests of each case before moving to the next and repeating 
until enough repetitions had been run. This was done in order to minimize any 
possible effect of varing environmental conditions on the test node over time.

\subsection{Analysis}

Finally, analysis was performed using the Jupyter notebook, as listed in the 
artifact appendix to this paper. The analysis is presented in section 
\ref{results}.

\section{Results} \label{results}

\subsection{Pilot Study}

We performed a pilot study to determine how many repetitions of each experiment 
were required for statistically significant results. For the D problem size, 
this indicated about 200 runs were required.

For the E problem size, analysis indicated about 106 to 108 repetitons were 
required, except in the case of HPCPerfStats, where 20 repetitions had already 
achieved statistical significance.

\subsection{Full Study}

Our full study showed the following results, organized by test case:

\begin{table}[h!]
\centering
\begin{tabular}{| c | c | c | c | c |} 
    \hline
    Test case & Overhead & test cases run & p (T-test) & p (TOST) \\ 
    \hline
    EP, size D & 0.38\% & 200 & 0.025 & 0.251 \\ 
    \hline
    EP, size E & 2.03\% & 20 & 0.007 & 0.981 \\ 
    \hline
\end{tabular}
\caption{Summary of results for HPCPerfStats}
\label{table:results_HPCPerfStats}
\end{table}

These results indicate a statistically significant impact of running 
HPCPerfStats in all cases.

\begin{table}[h!]
\centering
\begin{tabular}{| c | c | c | c | c |} 
    \hline
    Test case & Overhead & test cases run & p (T-test) & p (TOST) \\ 
    \hline
    EP, size D & 0.11\% & 200 & 0.481 & 0.009 \\ 
    \hline
    EP, size E & 0.05\% & 110 & 0.859 & 0.047 \\ 
    \hline
\end{tabular}
\caption{Summary of results for LDMS}
\label{table:results_LDMS}
\end{table}

These results indicate that running LDMS is within our bound for being 
considered equivalent to no monitoring in all cases.

\begin{table}[h!]
\centering
\begin{tabular}{| c | c | c | c | c |} 
    \hline
    Test case & Overhead & test cases run & p (T-test) & p (TOST) \\ 
    \hline
    EP, size D & 0.07\% & 200 & 0.747 & 0.015 \\ 
    \hline
    EP, size E & 0.04\% & 110 & 0.869 & 0.036 \\ 
    \hline
\end{tabular}
\caption{Summary of results for BMC}
\label{table:results_BMC}
\end{table}

These results indicate that the impact of BMC-based monitoring is within our 
predetermined bound for being considered equivalent to no monitoring.

% E job size, 20 repeats
% HPCPerfStats: ep_test: mops_total: 20/20 experimenatl/base samples
%     t=2.8494311352065678 
%     p-value, different=0.0070355411310043655
%     p-value, equivalent=0.9809197482535233 (lower: 0.9809197482535233, upper: 0.0005214041045588922)
%     baseline:     mean: 2755.29, stddev: 51.02268436685783
%     experimental: mean: 2699.3334999999997, stddev: 68.73056730996767
%     difference: -2.030875152887726%
% HPCPerfStats: ep_test: mops_process: 20/20 experimenatl/base samples
%     t=2.846794666098394 
%     p-value, different=0.007083576476303162
%     p-value, equivalent=0.9808136916791611 (lower: 0.9808136916791611, upper: 0.0005255436329799833)
%     baseline:     mean: 49.201499999999996, stddev: 0.9118896588951979
%     experimental: mean: 48.202999999999996, stddev: 1.2271434308995837
%     difference: -2.0294096724693356%
% HPCPerfStats: ep_test: runtime: 20/20 experimenatl/base samples
%     t=-2.840043097524336 
%     p-value, different=0.007207986279967026
%     p-value, equivalent=0.9816381631969431 (lower: 0.0005775261923905083, upper: 0.9816381631969431)
%     baseline:     mean: 798.3884999999999, stddev: 15.06863904770435
%     experimental: mean: 815.1865, stddev: 20.919506035038214
%     difference: 2.103988221273242%

% D job size, 200 repeats
% dict_keys(['ep_test'])
% HPCPerfStats: ep_test: mops_total: 200/200 experimenatl/base samples
%     t=2.243478359389888 
%     p-value, different=0.02541532833000295
%     p-value, equivalent=0.25121730156220407 (lower: 0.25121730156220407, upper: 1.9699379489396529e-07)
%     baseline:     mean: 2755.66725, stddev: 43.852086483284914
%     experimental: mean: 2745.0620999999996, stddev: 50.23690684735676
%     difference: -0.38484871495280576%
% HPCPerfStats: ep_test: mops_process: 200/200 experimenatl/base samples
%     t=2.245839615935669 
%     p-value, different=0.025262271922677357
%     p-value, equivalent=0.251831973829846 (lower: 0.251831973829846, upper: 1.9426030236667705e-07)
%     baseline:     mean: 49.208600000000004, stddev: 0.7827285864206059
%     experimental: mean: 49.01905000000001, stddev: 0.8971630830010785
%     difference: -0.3851968964774389%
% HPCPerfStats: ep_test: runtime: 200/200 experimenatl/base samples
%     t=-2.2278936966005314 
%     p-value, different=0.02644586948681947
%     p-value, equivalent=0.2757877789813298 (lower: 3.342352867448112e-07, upper: 0.2757877789813298)
%     baseline:     mean: 49.88785, stddev: 0.8158228223701515
%     experimental: mean: 50.08465, stddev: 0.9419282761972908
%     difference: 0.39448482947251323%
% dict_keys(['ep_test'])
% LDMS: ep_test: mops_total: 200/200 experimenatl/base samples
%     t=0.7037729293191733 
%     p-value, different=0.4819857200821027
%     p-value, equivalent=0.009272771511990439 (lower: 0.009272771511990439, upper: 9.333866148490351e-05)
%     baseline:     mean: 2755.66725, stddev: 43.852086483284914
%     experimental: mean: 2752.5066500000003, stddev: 45.722306462792325
%     difference: -0.11469454448826157%
% LDMS: ep_test: mops_process: 200/200 experimenatl/base samples
%     t=0.7064648673167511 
%     p-value, different=0.48031245588329574
%     p-value, equivalent=0.009332063026148467 (lower: 0.009332063026148467, upper: 9.226549216093664e-05)
%     baseline:     mean: 49.208600000000004, stddev: 0.7827285864206059
%     experimental: mean: 49.15195, stddev: 0.816657025623364
%     difference: -0.11512215344473273%
% LDMS: ep_test: runtime: 200/200 experimenatl/base samples
%     t=-0.7030605329827896 
%     p-value, different=0.48242906576851435
%     p-value, equivalent=0.011283129416094936 (lower: 0.00012497462320053358, upper: 0.011283129416094936)
%     baseline:     mean: 49.88785, stddev: 0.8158228223701515
%     experimental: mean: 49.946450000000006, stddev: 0.846714767498477
%     difference: 0.11746347056448721%
% dict_keys(['ep_test'])
% BMC: ep_test: mops_total: 200/200 experimenatl/base samples
%     t=0.3227556156649431 
%     p-value, different=0.7470498167796232
%     p-value, equivalent=0.01574031375547661 (lower: 0.01574031375547661, upper: 0.0026465682426937685)
%     baseline:     mean: 2755.66725, stddev: 43.852086483284914
%     experimental: mean: 2753.87505, stddev: 64.90682870852264
%     difference: -0.06503687990630459%
% BMC: ep_test: mops_process: 200/200 experimenatl/base samples
%     t=0.3242838477987278 
%     p-value, different=0.7458935056013593
%     p-value, equivalent=0.015784206755164773 (lower: 0.015784206755164773, upper: 0.002631036445601578)
%     baseline:     mean: 49.208600000000004, stddev: 0.7827285864206059
%     experimental: mean: 49.17645, stddev: 1.1590150549065357
%     difference: -0.06533410826563132%
% BMC: ep_test: runtime: 200/200 experimenatl/base samples
%     t=-0.4589878806187285 
%     p-value, different=0.6464934306588798
%     p-value, equivalent=0.03462691123631684 (lower: 0.003213505192534917, upper: 0.03462691123631684)
%     baseline:     mean: 49.88785, stddev: 0.8158228223701515
%     experimental: mean: 49.938050000000004, stddev: 1.3095333892268648
%     difference: 0.10062570345285235%

% LDMS: ep_test: mops_total: 110/110 experimenatl/base samples
%     t=0.1784546104359264 
%     p-value, different=0.8585316921714162
%     p-value, equivalent=0.046549918554651876 (lower: 0.046549918554651876, upper: 0.02109952548472329)
%     baseline:     mean: 2751.204454545455, stddev: 53.001269497426605
%     experimental: mean: 2749.8882727272726, stddev: 55.85835576558541
%     difference: -0.04784020380629766%
% LDMS: ep_test: mops_process: 110/110 experimenatl/base samples
%     t=0.17880283484251297 
%     p-value, different=0.8582585799291564
%     p-value, equivalent=0.046556294207266966 (lower: 0.046556294207266966, upper: 0.02106809982919118)
%     baseline:     mean: 49.12845454545454, stddev: 0.9464865041189647
%     experimental: mean: 49.10490909090909, stddev: 0.997142959983486
%     difference: -0.04792630821241096%
% LDMS: ep_test: runtime: 110/110 experimenatl/base samples
%     t=-0.18868237654473533 
%     p-value, different=0.8505172983879924
%     p-value, equivalent=0.05330487033806139 (lower: 0.02349705436925388, upper: 0.05330487033806139)
%     baseline:     mean: 799.6007272727271, stddev: 15.881313073093798
%     experimental: mean: 800.0177272727273, stddev: 16.73861711780662
%     difference: 0.05215102810404445%
% dict_keys(['ep_test'])
% BMC: ep_test: mops_total: 110/110 experimenatl/base samples
%     t=0.16502836713966312 
%     p-value, different=0.8690745808466529
%     p-value, equivalent=0.03629205218313655 (lower: 0.03629205218313655, upper: 0.016970202911894404)
%     baseline:     mean: 2751.204454545455, stddev: 53.001269497426605
%     experimental: mean: 2750.051636363637, stddev: 50.09867173937437
%     difference: -0.041902308638435415%
% BMC: ep_test: mops_process: 110/110 experimenatl/base samples
%     t=0.15811770106618397 
%     p-value, different=0.8745104201208962
%     p-value, equivalent=0.03577810518395838 (lower: 0.03577810518395838, upper: 0.017272519392580583)
%     baseline:     mean: 49.12845454545454, stddev: 0.9464865041189647
%     experimental: mean: 49.10872727272727, stddev: 0.8948977687584911
%     difference: -0.04015447444823269%
% BMC: ep_test: runtime: 110/110 experimenatl/base samples
%     t=-0.144223233621064 
%     p-value, different=0.8854574331616234
%     p-value, equivalent=0.03908971297308585 (lower: 0.020383947288111097, upper: 0.03908971297308585)
%     baseline:     mean: 799.6007272727271, stddev: 15.881313073093798
%     experimental: mean: 799.9019999999999, stddev: 14.947272223020855
%     difference: 0.037677895604271906%

\section{Discussion} \label{discussion}

These results lead to some interesting conclusions. We found that running 
HPCPerfStats had a statistically significant overhead on the system in all tested 
cases, though the degree of impact varied (0.38\% in size D, 2.0\% in size E). 
In contrast, the overhead for both LDMS and BMC-based monitoring was within our 
predefined bound for equivalence.

It is unclear why HPCPerfStats has a significant impact while LDMS does not. 
One theory is the number of modules performing collection - 31 for HPCPerfStats 
(although some hardware- and software-specific modules have no data to collect) 
versus LDMS's 6. Perhaps one (or more) of these modules have significant 
system impact. Benchmarking individual modules of monitoring systems is beyond 
the scope of this paper. However, the framework proposed within this paper 
could be used for such benchmarking.

Additionally, it is interesting that the measured impact of HPCPerfStats seems 
to vary by benchmark size. We measured 0.38\% when using EP problem size D and 
2.03\% with EP problem size E. Problem size D took approximately 50 seconds to 
run, while problem size E took approximately 800 seconds (13 minutes). We 
therorize that startup or teardown costs within the EP kernel may be masking 
the impact with problem size D. However, examining this theory is beyond the 
scope of the current work.

We confirm LDMS's statement that it runs without significant 
overhead on the system. Although from a 
statistcal point of view, we can only say that the overhead is less than 0.5\%,
the measured value of 0.11\% is significantly smaller and further statistical 
tests could more precisely define the overhead.

Finally, we confirm the commonly-accpeted theory that monitoring via the BMC 
does not have any meaningful overhead (with meaninful overhead defined as 
0.5\% or greater), as the two one-sided T-tests indicate the overhead is less 
than this threshold with p=0.036 for EP with problem size E.

\section{Future Work} \label{futurework}

This work focused on building and proving a methodology. It was 
specifically limited to single-node runs of CPU-focused benchmarks.

Future work will extend beyond these limits, examining how this methodology 
works on multi-node benchmarks. Further, it will extend to
other types of benchmarks - for example, GPU or tensor-hardware focused.

Additionally, future work will examine the applications of
out-of-band monitoring data. As we have now shown that collecting this data has 
no measurable cost of the performance of HPC systems, we wish to explore how
it can be used to better monitor running systems and jobs. For instance, can it 
be used to detect stalled jobs? Can it be used to determine what kinds of jobs
are running on a system? Can it be used for early detection of hardware 
failures?

\section{Conclusion} \label{conclusion}

This work contributes a methodology for measuring the impact of various 
monitoring systems on the performance of HPC systems as well as providing 
various levels of reporting in order to encourage reproducability. Further, it 
contributes a refernce implementation of this methodology in order to assist 
future usage by others. It then 
applies this methodology to measure overhead for HPCPerfStats, LDMS and 
BMC-based monitoring. Findings include a 2.0\% overhead for HPCPerfStats, 
whereas LDMS and BMC-based monitoring are both demonstrated to have an overhead 
of less than 0.5\%. This provides the first - to our knowledge - meaningful 
confirmation of the common-acepted theory that BMC-based monitoring has no 
impact on the running system.

\bibliography{Paper}{}

\end{document}
